\documentclass[12pt,a4paper]{article}
\usepackage[ngerman]{babel}
\usepackage[utf8]{inputenc}
\usepackage[T1]{fontenc}
\usepackage{geometry}
\usepackage{amsmath}
\usepackage{graphicx}
\usepackage{hyperref}
\usepackage{enumitem}
\usepackage{csquotes}

\geometry{margin=2.5cm}
\hypersetup{hidelinks}

\title{Künstliche Intelligenz als Herausforderung für Polizei und Justiz in Deutschland\\
	\large Eine Analyse rechtlicher, ethischer und praktischer Problemfelder}
\author{}
\date{\today}

\begin{document}
	
	\maketitle
	
	\begin{abstract}
		Die zunehmende Nutzung Künstlicher Intelligenz (KI) durch Polizei und Justiz in Deutschland eröffnet ein breites Spektrum neuer Herausforderungen. Die vorliegende Analyse systematisiert diese Problemfelder, von rechtlichen Grundlagendefiziten über das Risiko algorithmischer Diskriminierung bis hin zu konkreten Gefahren von Fehlurteilen durch KI-Halluzinationen und verzerrte Entscheidungsgrundlagen. Ein besonderer Fokus liegt auf dem neu entstehenden strukturellen Problem: dem strategischen Wettbewerbsvorteil privater Kanzleien gegenüber öffentlichen Institutionen, der die Fairness des Rechtsstaats gefährdet.
	\end{abstract}
	
	\section{Einleitung}
	
	Die Integration Künstlicher Intelligenz in die polizeiliche und justizielle Arbeit schreitet in Deutschland rasch voran. Dieser technologische Wandel bringt jedoch tiefgreifende Herausforderungen mit sich. Auf dem 3. IPA-Fachforum im Mai 2025 in Stuttgart brachte der IPA-Präsident die zentrale Fragestellung auf den Punkt: \enquote{Wie viel KI braucht unser Rechtsstaat – und wie viel hält er aus?} \cite{IPA2025}
	
	\section{Problemfelder beim polizeilichen Einsatz von KI}
	
	\subsection{Rechtliche Defizite und mangelnde Transparenz}
	
	Viele KI-Systeme im polizeilichen Kontext werden ohne ausreichende Rechtsgrundlage, Transparenz und Risikoabschätzung entwickelt, getestet oder genutzt \cite{Peteranderl2025}. Diese grundlegende Kritik wird durch die jüngste Debatte um die \enquote{Lex Palantir} bestätigt, bei der Bürgerrechtsgruppen einen \enquote{unverhältnismäßigen Eingriff} durch automatisierte Datenanalysen beklagen \cite{FR2026}.
	
	Das Bundesverfassungsgericht hat in seiner \enquote{Hessendata}-Entscheidung vom Februar 2023 verfassungsrechtliche Leitplanken für automatisierte Datenanalysen formuliert und die neuartige Eingriffsintensität dieser Technologie betont \cite{VerfBlog2026}. Die Konferenz der unabhängigen Datenschutzaufsichtsbehörden (DSK) forderte im September 2025 eine verfassungskonforme Ausgestaltung dieser Verfahren und betonte die Wahrung digitaler Souveränität \cite{DSK2025}.
	
	\subsection{Risiko algorithmischer Diskriminierung}
	
	KI-Systeme können strukturelle Vorurteile aus Trainingsdaten übernehmen und systematisch gegen bestimmte Gruppen diskriminieren. Die Forschungsstelle der Antidiskriminierungsstelle des Bundes warnt, dass der Einsatz von KI in der Polizei \enquote{sehr riskant} sei, da Algorithmen mitentscheiden könnten, welche Stadtteile verstärkt überwacht werden \cite{Antidisk2025}. \cite{taz2026} weisen darauf hin, dass Algorithmen Vorurteile wie Rassismus reproduzieren und Diskriminierung sogar verstärken können.
	
	Zivilgesellschaftliche Organisationen kritisieren insbesondere, dass automatisierte Auswertungen sich nicht auf Tatverdächtige beschränken, sondern auch Unbeteiligte erfassen und marginalisierte Gruppen dadurch besonders belasten können \cite{Netzpolitik2025}.
	
	\subsection{Ausufernde Überwachung und Grundrechtseingriffe}
	
	Der Einsatz KI-gestützter Systeme wie Gesichtserkennung in Echtzeit ermöglicht eine flächendeckende Überwachung des öffentlichen Raums. Im Frankfurter Bahnhofsviertel wurde im Juli 2025 ein entsprechendes Pilotprojekt gestartet \cite{TOnline2025}. Der Deutsche Anwaltverein warnt, dass biometrische Gesichtserkennung einen \enquote{schweren Eingriff in das Recht auf informationelle Selbstbestimmung} darstelle \cite{DAV2025}.
	
	Der Deutschlandfunk berichtet von Plänen des Bundesinnenministeriums, die Befugnisse für Gesichtserkennung massiv auszubauen, einschließlich des Abgleichs mit Bildern aus dem Internet \cite{DLF2025}. Ein Gutachten von Algorithm Watch zeigt zudem, dass polizeiliche Foto-Fahndung im Internet nur mit illegalen Datenbanken funktionieren kann \cite{taz2025}.
	
	Die \enquote{Lex Palantir} wird von Amnesty International als \enquote{rote Linie} kritisiert, die die Infrastruktur zur flächendeckenden Überwachung der Bevölkerung schaffe \cite{Amnesty2026}.
	
	\subsection{Das Blackbox-Problem und fehlende Wirksamkeitsbelege}
	
	Die oft undurchschaubaren Entscheidungswege von KI-Systemen können fatale Folgen haben. Der Deutsche Anwaltverein bemängelt, dass die Funktionsweise der Algorithmen völlig intransparent sei \cite{DAV2025}. \cite{Peteranderl2025} betonen, dass Entwicklung und Einsatz häufig ohne unabhängige wissenschaftliche Begleitung sowie Kontrollmechanismen erfolgen – obwohl die Systeme erhebliche Risiken bergen.
	
	Besonders problematisch ist, dass es für prädiktive Polizeiarbeit (\enquote{Predictive Policing}) in Deutschland bislang keinen wissenschaftlichen Nachweis für deren Wirksamkeit gibt \cite{Peteranderl2025}. \cite{CILIP2025} bestätigen diese Einschätzung: Es existiere kein wissenschaftlicher Beleg für den Einfluss von Predictive Policing auf eine sinkende Kriminalitätsrate.
	
	\subsection{Abhängigkeit von US-Tech-Konzernen}
	
	Die Nutzung der Analysesoftware \enquote{Gotham} des US-Unternehmens Palantir führt zu einer problematischen technologischen Abhängigkeit. In Baden-Württemberg löste der Alleingang von Innenminister Strobl beim Vertragsabschluss mit Palantir im März 2025 einen Koalitionskonflikt aus \cite{FR2026}. Kritiker befürchten Hintertüren für US-Geheimdienste \cite{Focus2025}, obwohl Palantir dies zurückweist \cite{Spiegel2025}. In Hessen wird die Software bereits seit 2017 eingesetzt, musste aber nach einem Urteil des Bundesverfassungsgerichts angepasst werden \cite{FR2026}.
	
	Einige Bundesländer wie Niedersachsen und Bremen setzen bewusst nicht auf Palantir \cite{Goslarsche2025}. In Baden-Württemberg strebt die Koalition eine europäische Alternative bis 2030 an \cite{FR2026}.
	
	\section{Problemfelder in der Justiz}
	
	\subsection{KI als Hilfsmittel, nicht als Richterersatz}
	
	Der Thüringer Richterbund zieht eine klare Grenze beim KI-Einsatz in der Justiz. Der Landesvorsitzende Holger Pröbstel erklärte gegenüber der Deutschen Presse-Agentur: \enquote{Was aus meiner Sicht gar nicht geht, ist, dass Sie sich von ChatGPT die Urteile schreiben lassen} \cite{SZ2025}. Die Findung eines Urteils müsse immer Aufgabe eines Richters sein: \enquote{Sonst können Sie die Justiz gleich ganz abschaffen} \cite{SZ2025}. KI dürfe für Gerichte und Staatsanwaltschaften immer nur ein Hilfsmittel sein, niemals den Richter ersetzen \cite{SZ2025}.
	
	\subsection{Gefährdung der richterlichen Unabhängigkeit}
	
	Eine tiefe inhaltliche Einbindung von KI in die Entscheidungsfindung kann die richterliche Unabhängigkeit untergraben. Aus der bisherigen Rechtsprechung ist abzuleiten, dass verpflichtende Vorgaben zur Nutzung von KI-Tools im richterlichen Kernbereich unzulässig sind \cite{Anwaltsblatt2025}.
	
	\subsection{Fehleranfälligkeit und mangelnde Standards}
	
	Skeptisch zeigt sich Pröbstel insbesondere gegenüber Überlegungen, mittels KI den Strafrahmen zu bestimmen. Die entscheidende Frage sei: \enquote{Welche Urteile stehen in dieser Datenbank und wie bewerten das die Algorithmen} \cite{SZ2025}. Da die Justiz auf diese Algorithmen keinen Zugriff habe, könne sie so nicht seriös arbeiten \cite{SZ2025}.
	
	Bisher fehlen zudem einheitliche Standards für den Umgang mit KI-generierten Beweismitteln. Die automatisierte Gesichtserkennung in der Strafverfolgung wird in Deutschland bereits seit rund zwei Jahrzehnten praktiziert – allerdings ohne taugliche Rechtsgrundlage \cite{Nomos2025}.
	
	\subsection{Fehlerhafte Verfahrenshandlungen durch KI-Halluzinationen}
	
	Eine der konkretesten Gefahren für Fehlurteile geht vom unkritischen Einsatz generativer KI in der anwaltlichen und richterlichen Arbeit aus. Große Sprachmodelle neigen dazu, plausible, aber frei erfundene Informationen zu produzieren. Eine im Juni 2025 veröffentlichte Datenbank des französischen Anwalts Damien Charlotin weist allein für das Jahr 2025 bereits über 80 vor Gericht aufgedeckte Fälle von KI-Halluzinationen nach \cite{Charlotin2025}.
	
	\subsubsection{Praxisbeispiel Amtsgericht Köln (2025)}
	Ein Rechtsanwalt reichte einen mit Hilfe von KI verfassten Schriftsatz ein, der nicht existierende Gerichtsurteile zitierte. Das Gericht rügte den Anwalt öffentlich und erkannte den Schriftsatz nicht als vollwertig an. Der Fall zeigt, wie ein Mangel an der anwaltlichen Sorgfaltspflicht das gesamte Verfahren beeinträchtigen kann \cite{LTO2025a}.
	
	Weltweit wurden inzwischen mehr als 1.200 gerichtliche Sanktionen gegen Rechtsanwälte wegen fehlerhafter KI-Nutzung verhängt, davon rund 800 allein in den USA \cite{NPR2025}. Einige Gerichte verlangen mittlerweile die Offenlegung der KI-Nutzung in Schriftsätzen.
	
	\subsubsection{Risiken für die Urteilsfindung}
	Wenn Richter Entscheidungen auf KI-Empfehlungen stützen, ohne die Algorithmen zu verstehen, entsteht eine \enquote{Blackbox-Justiz}. Studien aus dem Jahr 2025 zeigen, dass bis zu 15 Prozent der automatisierten Fallanalysen durch KI-Systeme fehlerhaft sind \cite{LTO2025b}. Eine groß angelegte Untersuchung der Europäischen Rundfunkunion stellte zudem eine Fehlerquote von bis zu 40 Prozent bei Antworten gängiger Chatbots fest \cite{heise2025}.
	
	\subsection{Manipulierte oder unzuverlässige Beweise als Grundlage von Fehlurteilen}
	
	Die zunehmende Qualität von Deepfakes ermöglicht die Manipulation audiovisueller Beweismittel. Im November 2025 reichten Kläger in einem US-amerikanischen Zivilprozess ein Video einer angeblichen Zeugin als Beweis ein, das sich später als Deepfake entpuppte. Die Richterin wies die Klage ab, nachdem sie technische Artefakte identifiziert hatte \cite{LTO2025c}. Auch wenn es sich um einen US-Fall handelt, verdeutlicht er das erhebliche Gefahrenpotenzial für deutsche Gerichte.
	
	Eine weitere Beweisrisikoquelle liegt in der unzulässigen Nutzung von KI durch gerichtlich bestellte Sachverständige. Das Landgericht Darmstadt stellte am 10. November 2025 klar, dass ein Sachverständigengutachten, das umfangreich mit KI erstellt wurde, als Beweismittel unzulässig ist. Das Gericht monierte, dass der Gutachter persönliche Untersuchungen unterlassen und sich statt auf die Akten auf KI-generierte Fallzusammenfassungen gestützt hatte, die nicht mit den Tatsachen übereinstimmten. Der Gutachter erhielt keinerlei Vergütung \cite{LG Darmstadt2025}. Solche mangelhaften Gutachten können im Zweifel zu Fehlurteilen beitragen.
	
	\subsection{Verzerrte Entscheidungsgrundlagen durch algorithmische Verzerrungen}
	
	Fehlurteile können auch durch inhärente Verzerrungen in den Trainingsdaten von KI-Systemen entstehen. Wenn ein Algorithmus aus historischen, möglicherweise diskriminierenden Gerichtsurteilen lernt, wird er diese Fehler wiederholen und potenziell verstärken. Ein Forschungsprojekt an der Universität Konstanz untersucht derartige \enquote{algorithmic biases} und ihre Auswirkungen auf die Strafjustiz \cite{Uni Konstanz2025}. Diese systematischen Verzerrungen sind besonders tückisch, da sie nicht auf einen einzelnen menschlichen Fehler zurückgehen, sondern flächendeckend in die Strafzumessung einfließen können.
	
	\subsection{Wettbewerbsvorteil privater Kanzleien durch KI}
	
	Über die unmittelbaren Gefahren von Fehlurteilen hinaus zeichnet sich ein strukturelles Problem ab: Private Kanzleien nutzen KI-Technologie strategisch und oft schneller als öffentliche Institutionen, was zu einem massiven \enquote{Ungleichgewicht der Waffen} (asymmetry of arms) führen kann.
	
	\subsubsection{Strategische Frühphase und unternehmerischer Druck}
	Kanzleien erkennen den Wettbewerbsvorteil durch KI und investieren gezielt. Eine im Oktober 2025 veröffentlichte Benchmark-Studie des Unternehmens Vals AI, die auf einem Datensatz von über 200 juristischen Fragestellungen aus acht US-Kanzleien basierte, zeigte, dass vier getestete KI-Produkte bei den meisten Aufgaben besser abschnitten als die menschliche Vergleichsgruppe. Die durchschnittliche Genauigkeit der KI-Produkte lag bei 80\%, während die der Jurist:innen 71\% erreichte. Besonders bei der Überprüfbarkeit von Quellen (\enquote{Belegbarkeit}) lagen spezialisierte Legal-KI-Produkte mit 76\% deutlich vor menschlichen Jurist:innen (68\%). Noch deutlicher war der Unterschied bei der Mandantenfreundlichkeit und Verständlichkeit (\enquote{Angemessenheit}) – hier erreichten die Jurist:innen nur 60\%, während die KI-Produkte auf 70\% kamen \cite{ValAI2025}.
	
	Prof. Dr. Frauke Rostalski, Direktorin des Instituts für Strafrecht und Strafprozessrecht an der Universität zu Köln, sieht darin ein grundlegendes Problem für die Fairness des Rechtsstaats. Sie betont, dass KI-Systeme für Privatpersonen vor allem eine unterstützende Rolle im Bereich der rechtlichen Selbsthilfe spielen können. Das daraus entstehende Machtgefälle könnte Bürger gegenüber einem übermächtigen Staat jedoch weiter benachteiligen, wenn der Staat hier nicht Schritt hält \cite{Rostalski2024}.
	
	\subsubsection{Strukturelle Ungleichgewichte in der Prozessführung}
	Dieser Effizienzvorteil wird in Massenverfahren besonders deutlich. Der Einsatz von KI bei der Dokumentenanalyse (Technology Assisted Review) ermöglicht es Kanzleien, riesige Datenmengen – tausende E-Mails, Verträge, Gutachten – extrem schnell zu durchsuchen und prozessrelevante Informationen zu extrahieren. Es ist zudem zunehmend üblich, KI-Tools wie \enquote{Harvey} als \enquote{Gedankenpartner} für die Entwicklung einer Prozessstrategie zu nutzen \cite{Hengeler2025}. Eine KI kann tausende relevante Urteile analysieren, Argumentationsmuster erkennen und dem gegnerischen Anwalt, der sich manuell durch die Akten kämpfen muss, einen entscheidenden Schritt voraus sein.
	
	\subsubsection{Konkrete Machtasymmetrie: Fallbeispiel}
	Ein konkretes Fallbeispiel verdeutlicht diese Asymmetrie: Eine Partei verfügt über Tausende von Schadensberichten zu einzelnen beschädigten Gegenständen. Das Unternehmen nutzte eine KI-Plattform, um Beweise vorzubereiten, indem es die Dokumente hochlud und die KI Fragen dazu stellen ließ. Die Gegenseite (oder das Gericht) müsste all diese Unterlagen manuell sichten, um die Beweise zu prüfen. Schon hier offenbaren sich die inhärenten Risiken: Die KI war nicht in der Lage, alle relevanten Informationen zu filtern \cite{Luther2026}. Während ein geschulter Mensch zwar in der Lage wäre, dies korrekt zu tun, würde er dafür unverhältnismäßig viel Zeit benötigen – ein Nachteil, der in der Praxis oft zu prozessualen Zugeständnissen oder Vergleichsverhandlungen führt.
	
	\subsubsection{Die Notwendigkeit eines \enquote{Sovereign AI}}
	Prof. Steinrötter von der Universität Potsdam weist darauf hin, dass die öffentliche Hand angesichts dieser Übermacht nicht nur die technologischen, sondern auch die rechtlichen Rahmenbedingungen schaffen muss. Er betont: \enquote{Letztlich müssen alle relevanten Gesetze, die dazugehörige Rechtsprechung und die Verwaltungspraxis in das System eingespeist werden.} Der Staat müsse insbesondere auf \enquote{Sovereign AI} – also souveräne, europäische KI-Lösungen – setzen, um nicht von US-Konzernen abhängig zu werden \cite{Steinroetter2026}.
	
	\subsubsection{Reputations- und Haftungsrisiken als Bremsfaktor für Kanzleien}
	Die Kehrseite des Wettbewerbsvorteils ist das immense Risiko bei unsachgemäßer Nutzung. Die aggressive Nutzung von KI kann für Kanzleien nach hinten losgehen, was ihre Zurückhaltung erklären könnte. Das Landgericht Frankfurt am Main rügte im September 2025 einen Anwalt scharf, weil dieser in einem Schriftsatz frei erfundene BGH-Zitate aus einem KI-Tool verwendet hatte (Az. 2-13 S 56/24). Die Richter sprachen von einer \enquote{kompletten Fälschung} \cite{LG Frankfurt2025}. Das Landgericht Darmstadt kürzte in einem Fall das Honorar eines Sachverständigen auf 0,00 Euro, weil er ein maßgeblich von einer KI erstelltes Gutachten nicht deklariert hatte (Az. 19 O 527/16) \cite{Kramarz2026}. In einem anderen Fall wertete das Kammergericht Berlin automatisierte Schriftsätze ohne Bezug zum konkreten Tatvorwurf als dysfunktionale Prozessführung \cite{KG Berlin2025}.
	
	\subsection{Reaktionen der Gerichte und Grenzen des richterlichen Ersatzes}
	
	Der österreichische Oberste Gerichtshof (OGH) wies im Oktober 2025 eine offensichtlich mit KI erstellte Nichtigkeitsbeschwerde zurück, die sich auf erfundene angebliche höchstgerichtliche Entscheidungen stützte. Das Gericht stellte klar, dass dieses Vorbringen dem erforderlichen Argumentationsniveau nicht genüge \cite{OGH2025}.
	
	Der Thüringer Richterbund betonte im August 2025: \enquote{Was aus meiner Sicht gar nicht geht, ist, dass Sie sich von ChatGPT die Urteile schreiben lassen} \cite{Richterbund2025}. Die Findung eines Urteils müsse immer einem Richter vorbehalten bleiben. KI dürfe für Gerichte immer nur ein Hilfsmittel sein, aber \enquote{nie den Richter ersetzen} \cite{Richterbund2025}.
	
	\section{Die Rolle der deutschen und europäischen IT-Industrie und der Open-Source-Community}
	
	Angesichts der skizzierten Herausforderungen stellt sich zwangsläufig die Frage, inwieweit die öffentliche Hand die notwendige technologische Souveränität überhaupt erreichen kann, ohne sich noch tiefer in Abhängigkeiten von großen US-Tech-Konzernen zu begeben. Dies erfordert einen differenzierten Blick auf die Strategien der Bundesregierung, der europäischen IT-Wirtschaft und der Open-Source-Community.
	
	\subsection{Die Open-Source-Strategie des Bundes: Das Zentrum für Digitale Souveränität (ZenDiS)}
	
	Um die öffentliche Verwaltung aus der Abhängigkeit von wenigen großen Tech-Konzernen zu befreien, wurde im Dezember 2022 das \textbf{Zentrum für Digitale Souveränität der Öffentlichen Verwaltung (ZenDiS)} gegründet \cite{ZenDiS2025}. Das ZenDiS ist eine GmbH im Besitz des Bundes und verfolgt das Ziel, Bund, Ländern und Kommunen ein umfassendes Souveränitätspaket aus Plattform, Produkten und Beratung zur Verfügung zu stellen \cite{ZenDiS2025}. Kern der Strategie ist das Prinzip \enquote{Public Money, Public Code}, das im Koalitionsvertrag 2021 verankert wurde und besagt, dass mit öffentlichen Geldern entwickelte Software auch der Allgemeinheit als Open Source zur Verfügung stehen soll \cite{SovereignTechFund2025}.
	
	Das ZenDiS hat zwei zentrale Plattformen geschaffen:
	\begin{itemize}
		\item \textbf{openCode:} Eine offene Plattform (basierend auf GitLab) zur Entwicklung, Bereitstellung und Nachnutzung von Open-Source-Software für die öffentliche Verwaltung. Sie dient als zentrale Anlaufstelle für die kollaborative Softwareentwicklung von Bund, Ländern und Kommunen. Mittlerweile basieren über 4.000 Projekte auf openCode \cite{Heise2025a}.
		\item \textbf{openDesk:} Eine souveräne Open-Source-Arbeitsplatzlösung, die als vollwertige Alternative zu Microsoft 365 und Teams konzipiert ist. openDesk integriert etablierte Open-Source-Tools wie Open-Xchange (Groupware), OpenProject (Projektmanagement), Nextcloud (Dateisynchronisation), XWiki (Wissensmanagement), Matrix und Jitsi (Kommunikation) sowie Collabora (Web-Office) in einer einheitlichen Plattform \cite{AllThingsOpen2025}. Die Verfügbarkeit als Source-Code, Helm-Charts und bald auch als SaaS-Angebot macht es zu einem echten Werkzeug für die digitale Souveränität \cite{AllThingsOpen2025}.
	\end{itemize}
	
	\subsection{Die Open-Source-Community als Fundament: Der Sovereign Tech Fund}
	
	Neben der Bereitstellung fertiger Lösungen hat die Bundesregierung ein international vielbeachtetes Instrument zur gezielten Stärkung der Open-Source-Infrastruktur geschaffen: den \textbf{Sovereign Tech Fund (STF)}. Der STF ist ein Förderprogramm des Bundesministeriums für Wirtschaft und Klimaschutz (BMWK), das darauf abzielt, die Sicherheit, Wartung und Resilienz kritischer Open-Source-Basistechnologien zu stärken \cite{SovereignTechFund2025}. Seit seiner Gründung im Jahr 2022 hat der STF über 24,6 Millionen Euro bereitgestellt, um mehr als 60 Open-Source-Projekte weltweit zu unterstützen \cite{SovereignTechFund2025}. Die Mindestförderung pro Projekt beträgt 50.000 Euro, in der Praxis wurden bereits Investitionen von bis zu einer Million Euro getätigt \cite{SovereignTechFund2025}.
	
	Der STF adressiert damit die strukturelle Herausforderung, dass viele grundlegende Open-Source-Komponenten – wie Bibliotheken, Protokolle und Entwicklungswerkzeuge – aufgrund ihrer Allgegenwart in der digitalen Infrastruktur oft vernachlässigt werden, obwohl sie das Rückgrat der digitalen Wirtschaft und Verwaltung bilden. Mit Programmen wie dem \enquote{Sovereign Tech Fellowship}, das Maintainer*innen von wichtigen Open-Source-Komponenten für ihre Arbeit bezahlt, verfolgt der STF einen ganzheitlichen Ansatz zur Sicherung des Open-Source-Ökosystems im öffentlichen Interesse \cite{STFNews2025}.
	
	\subsection{Deutsche IT-Wirtschaft: Zwischen Eigenentwicklung und Abhängigkeit}
	
	Die deutsche IT-Wirtschaft ist in diesem Ökosystem auf verschiedene Weisen aktiv. Einerseits gibt es vielversprechende Eigenentwicklungen, die bewusst auf Open Source setzen, um Unabhängigkeit zu gewährleisten.
	
	Ein herausragendes Beispiel ist der KI-Assistent \textbf{F13}, der vom Innovationslabor Baden-Württemberg entwickelt und im Juli 2025 als Open-Source-Software (MPL 2.0 Lizenz) auf openCode veröffentlicht wurde \cite{Staatsanzeiger2025}. F13 ist eine modellagnostische, datensouveräne KI-Assistenz, die auf deutschen Servern betrieben wird und somit DSGVO-konform ist. Sie umfasst bereits Funktionen wie Chat, Dokumentenzusammenfassung und Recherche in verwaltungsrelevanten Quellen und wird bereits von der Polizei Baden-Württemberg genutzt \cite{Staatsanzeiger2025}. Das ZenDiS hat das Projekt bei der Lizenzprüfung und der Community-Strategie beratend unterstützt \cite{ZenDiS2025b}. Diese Entwicklung zeigt, dass eine datenschutzkonforme und souveräne KI-Infrastruktur für die öffentliche Hand technisch möglich ist – wenn der politische Wille und die entsprechenden Ressourcen vorhanden sind.
	
	Andererseits wird die anhaltende Abhängigkeit von proprietären US-Lösungen besonders deutlich im Bereich der Polizei-Datenanalyse. Die \textbf{Palantir}-Software ist in den Polizeibehörden von Bayern, Hessen, Baden-Württemberg und Nordrhein-Westfalen im Einsatz \cite{Tagesschau2025}. Politisch ist man sich einig, dass eine souveräne Alternative geschaffen werden muss. So bekräftigte die Innenministerkonferenz im Juni 2025 das Ziel, eine \enquote{Verfahrensübergreifende Recherche- und Analyseplattform} zu schaffen, die frei von außereuropäischen Einflüssen ist \cite{Tagesschau2025}.
	
	Dennoch scheiterten vielversprechende Eigeninitiativen. So schloss sich 2023 ein Konsortium aus sieben deutschen Unternehmen (Start-Ups, Mittelständler und ein Großkonzern) unter dem Namen \textbf{NaSa} zusammen, um eine konkrete Palantir-Alternative zu entwickeln. Sachar Paulus, Professor für IT-Sicherheit an der Hochschule Mannheim, hielt eine solche Lösung für realistisch und schätzte, dass innerhalb von zwölf bis 18 Monaten ein Prototyp umsetzbar sei \cite{Tagesschau2025}. Trotz eines Zeitplans und der Vorstellung im Innenausschuss des Bundestags fand das Projekt keinen Auftraggeber – weder auf Landes- noch auf Bundesebene, sodass es scheiterte \cite{Tagesschau2025}. Dies zeigt ein strukturelles Problem auf: Während die Politik digitale Souveränität propagiert, scheitert die konkrete Umsetzung oft an fehlenden verbindlichen Aufträgen und der mangelnden Bereitschaft, langfristig in heimische Entwicklungen zu investieren.
	
	\subsection{Europäische Initiativen: Open Source in der Justiz}
	
	Auch im Justizbereich gibt es ermutigende Open-Source-Projekte, die den Weg für mehr Souveränität ebnen. So gewann das Bundesministerium der Justiz und für Verbraucherschutz mit dem Projekt \textbf{„Zugang zum Recht“} den Open-Source-Wettbewerb der Open Source Business Alliance (OSBA) in der Kategorie „Fachverfahren“ \cite{Heise2025c}. Die Anwendung bietet digitale Justizdienste an, allen voran die digitale Klage für Fluggastrechte, die es Bürgern ermöglicht, direkt eine digitale Klage ohne Anwalt einzureichen. Jährlich ließen sich so rund 790.000 Euro an Beratungshilfekosten einsparen \cite{Heise2025c}.
	
	Ein weiteres Leuchtturmprojekt ist \textbf{NeuRIS}, das neue Rechtsinformationssystem des Bundes. Der DigitalService entwickelt gemeinsam mit dem BMJV und dem Bundesamt für Justiz ein Portal, das Gesetze, Gerichtsentscheidungen und Verwaltungsvorschriften des Bundes zentral, barrierefrei und über eine offene API kostenlos zur Verfügung stellt \cite{NeuRIS2025}. Dies ermöglicht nicht nur einen modernen Zugang zum Recht, sondern schafft auch eine datensouveräne Grundlage für Legal-Tech-Anwendungen.
	
	\subsection{Ausblick: Open Source als Standard}
	
	Die Weichen für eine stärkere Open-Source-Nutzung sind gestellt. Die überarbeiteten EVB-IT-Verträge, die im Februar 2026 veröffentlicht wurden, machen Open Source zum Standard für die öffentliche Verwaltung und standardisieren die rechtskonforme Beschaffung von Open-Source-Software \cite{ZenDiS2025}. Auch im EU-weiten Kontext wird das Thema vorangetrieben: Im Oktober 2025 wurde das \enquote{EU Digital Commons Consortium} (DC-EDIC) offiziell gegründet, um digitale Gemeingüte in Europa zu fördern. Deutschland wird hier durch das Bundesministerium für Digitales und Staatsmodernisierung vertreten \cite{ZenDiS2025}.
	
	\section{Verbindende Querschnittsthemen}
	
	\subsection{Datenschutz und Grundrechte}
	
	Die automatisierte Analyse personenbezogener Daten durch KI stellt einen massiven Eingriff in die informationelle Selbstbestimmung dar. Die DSK betont, dass automatisierte Datenanalysen grundsätzlich alle Menschen betreffen können, ohne dass sie durch ihr Verhalten einen Anlass für polizeiliche Ermittlungen gegeben hätten \cite{DSK2025}.
	
	Besonders problematisch sind laut DSK Datenanalysen von Personen, die keinen Anlass für Ermittlungen gegeben haben, sowie Zweckänderungen \cite{Datenschutzticker2025}. Die Humanistische Union warnt, dass KI-gestützte \enquote{Superdatenbanken} im Widerspruch zur EU-KI-Verordnung stehen und weitreichende Profilbildung auch von Unbeteiligten ermöglichen \cite{HumanUnion2025}.
	
	\subsection{EU-KI-VO als neuer Rechtsrahmen}
	
	Die EU-KI-Verordnung (KI-VO) reguliert KI-Systeme nach einem risikobasierten Ansatz und unterwirft Hochrisiko-KI-Systeme besonders strengen Anforderungen \cite{Kriminalpolizei2026}. Biometrische Massenüberwachung in Echtzeit ist nach Art. 5 KI-VO grundsätzlich eine verbotene Praxis \cite{RAV2025}.
	
	Die Bundesregierung plant jedoch offenbar eine Auslagerung des biometrischen Abgleichs ins nicht-europäische Ausland, um diese Verbote zu umgehen \cite{Gruene2026}. Innenminister Dobrindt fordert sogar eine \enquote{Korrektur} der Leitlinien für Hochrisiko-Systeme \cite{Dobrindt2026}. Die KI-VO erstreckt sich nach überwiegender Auffassung auch auf präventiv polizeiliche Aufgabenwahrnehmung \cite{Kriminalpolizei2026}.
	
	\section{Fazit}
	
	KI für Polizei und Justiz stellt kein eindimensionales Problem dar, sondern eröffnet ein breites Spektrum neuer Herausforderungen. Die größten Hürden liegen im Spannungsfeld zwischen technologischen Möglichkeiten und rechtsstaatlichen Prinzipien: bei der Frage nach angemessenen rechtlichen Rahmenbedingungen, dem Schutz vor Diskriminierung, der Wahrung richterlicher Unabhängigkeit und der Sicherung informationeller Selbstbestimmung. Besonders alarmierend sind die konkreten Gefahren von Fehlurteilen durch KI-Halluzinationen, manipulierte Beweise und verzerrte Entscheidungsgrundlagen, die bereits jetzt zu nachweisbaren Verfahrensfehlern führen. Hinzu tritt ein neuartiges strukturelles Problem: Private Kanzleien nutzen KI strategisch zu ihrem Vorteil, was zu einem gefährlichen Ungleichgewicht im Rechtsstaat führen kann, wenn der Staat technologisch nicht Schritt hält.
	
	Die strategische Antwort des Bundes auf diese Herausforderung – ZenDiS, STF, openCode, openDesk, F13 und NeuRIS – zeigt, dass die technologischen und konzeptionellen Grundlagen für eine souveräne digitale Verwaltung vorhanden sind. Die entscheidende Frage ist nun, ob der politische Wille ausreicht, um diese Werkzeuge konsequent einzusetzen und die Entwicklung durch verbindliche Ausschreibungen und langfristige Investitionen zu fördern. Nur so kann Deutschland die drohende digitale Spaltung in eine \enquote{Zwei-Klassen-Justiz} verhindern und die Fairness des Rechtsstaats auch im Zeitalter der KI bewahren.
	
	\newpage
	\begin{thebibliography}{99}
		
		% ... (alle bisherigen Quellen wie zuvor)
		
		% Neue Quellen für Abschnitt 4
		\bibitem{ZenDiS2025} ZenDiS – Zentrum für Digitale Souveränität der Öffentlichen Verwaltung (2025): \emph{Startseite}. Abrufbar unter \url{www.zendis.de}.
		
		\bibitem{SovereignTechFund2025} Europäische Kommission (2025): \emph{Funding open source: case study on the Sovereign Tech Fund}. Open Source Observatory (OSOR), 19. Juni 2025. S. 1-53.
		
		\bibitem{STFNews2025} Sovereign Tech Agency (2025): \emph{Newsletter: Neue Technologie-Investitionen, Vorträge auf dem Open Source Summit und Sovereign Tech Fellowship Highlights}. 21. Oktober 2025.
		
		\bibitem{Heise2025a} Heise online (2025): \emph{Souveräne Verwaltung: Preisverleihung für offene Behördensoftware}. Bericht vom 15. Oktober 2025.
		
		\bibitem{AllThingsOpen2025} All Things Open (2025): \emph{ZenDIS, openDesk, and openCode: How Germany is transforming their public sector with open source}. 30. Januar 2025.
		
		\bibitem{Staatsanzeiger2025} Staatsanzeiger für Baden-Württemberg (2025): \emph{F13 wird Open Source: Baden-Württemberg bringt KI in die Verwaltung}. 23. Juli 2025.
		
		\bibitem{ZenDiS2025b} Zentrum Digitale Souveränität (2025): \emph{Beitrag auf social.bund.de zur Veröffentlichung von F13}. 24. Juli 2025.
		
		\bibitem{Tagesschau2025} Tagesschau (2025): \emph{Polizei-Datenanalysen – Könnte Palantirs Software ersetzt werden?} Bericht von L. Jeric und L. Hagen, 20. August 2025.
		
		\bibitem{Heise2025c} Heise online (2025): \emph{Souveräne Verwaltung: Preisverleihung für offene Behördensoftware}. Bericht vom 15. Oktober 2025.
		
		\bibitem{NeuRIS2025} DigitalService des Bundes (2025): \emph{NeuRIS: Service für Rechtsinformationen in der Testphase veröffentlicht}. Blogbeitrag vom 2. Juli 2025.
		
		\bibitem{OSBA2025} Open Source Business Alliance (2024): \emph{38 Maßnahmen für eine zukunftsfähige und Open-Source-getriebene Digitalisierung von Wirtschaft und Verwaltung}. Oktober 2024.
		
		\bibitem{Bitkom2025} Bitkom Research (2025): \emph{Open Source Monitor 2025}. Studie im Auftrag von Bitkom e. V.
		
		% Alle vorherigen Quellen hier einfügen...
		
	\end{thebibliography}
	
\end{document}